% =============================================================================
%  s01_intro.tex — Sección 1: Introducción / Planteamiento del problema
%  Sustituye \lipsum con tu contenido real.
% =============================================================================

\section{Introducción}

% ── Frame: Contexto ───────────────────────────────────────────────────────────
% NOTA \pause: revela bullets uno a uno durante la presentación en vivo,
%   pero genera páginas PDF duplicadas. Descomenta solo si presentas en proyector.
\begin{frame}{Contexto y Problemática}
  \begin{columns}[T]
    \begin{column}{0.58\textwidth}
      \begin{itemize}
        \item La \zmvm\ crece a una tasa anual de \textbf{1.4\%}, con el 80\% del suelo
              nuevo en zonas periféricas sin servicio de transporte masivo.
        % \pause
        \item La expansión horizontal eleva los \alert{costos de infraestructura}
              y reduce la accesibilidad para los estratos más vulnerables.
        % \pause
        \item El \termino{modelo de densificación compacta} emerge como alternativa
              con respaldo técnico y normativo creciente \citep{lore2024densidad}.
      \end{itemize}
    \end{column}
    \begin{column}{0.38\textwidth}
      \begin{block}{Dato clave}
        El 60\% de los viajes en la \zmvm\ se realizan en transporte informal,
        señal de un sistema de movilidad estructuralmente insuficiente.
      \end{block}
    \end{column}
  \end{columns}
\end{frame}

% ── Frame: Objetivos ─────────────────────────────────────────────────────────
\begin{frame}{Objetivos de la Presentación}
  \begin{block}{Objetivo General}
    Evaluar tres modelos de densificación urbana bajo criterios de equidad,
    eficiencia y viabilidad institucional en el contexto mexicano.
  \end{block}
  \vspace{0.4cm}
  \begin{columns}[T]
    \begin{column}{0.32\textwidth}
      \centering
      {\color{ColorPrincipal}\Large\textbf{01}}\\[0.1cm]
      {\small Analizar fundamentos teóricos}
    \end{column}
    \begin{column}{0.32\textwidth}
      \centering
      {\color{ColorPrincipal}\Large\textbf{02}}\\[0.1cm]
      {\small Identificar casos internacionales}
    \end{column}
    \begin{column}{0.32\textwidth}
      \centering
      {\color{ColorPrincipal}\Large\textbf{03}}\\[0.1cm]
      {\small Proponer lineamientos de política}
    \end{column}
  \end{columns}
\end{frame}
